% ============================================================================
% Section II: Literature Review and State of the Art
% ============================================================================

\section{Literature Review}
\label{sec:literature}

\subsection{Generation and Transmission Expansion Planning}

GTEP is a classical problem in power systems that has been studied
extensively since the 1970s.  Romero et al.~\cite{romero2002test} provided
canonical test systems and mathematical models---including transportation,
hybrid, DC power flow, and disjunctive formulations---that remain standard
benchmarks.  Hemmati et al.~\cite{hemmati2013comprehensive} presented a
comprehensive review covering generation expansion planning (GEP),
transmission expansion planning (TEP), and their integration.  More
recently, Mahdavi et al.~\cite{mahdavi2019tep_survey} classified TEP
literature into categories by mathematical model, solution method, and
uncertainty handling, identifying a trend toward stochastic,
multi-objective, and market-based formulations.

Lumbreras and Ramos~\cite{lumbreras2016progressive} surveyed modern
challenges in TEP including renewable integration, energy storage, and
the need for multi-stage stochastic models.  These challenges motivate
the development of tools that can handle joint generation, transmission,
and storage expansion under operational uncertainty.

\subsection{DC Optimal Power Flow}

The DC power flow approximation linearizes the AC power flow equations
by assuming lossless lines, flat voltage profiles, and small angle
differences~\cite{stott2009dc_opf}.  Despite its simplifications, DC OPF
is the standard formulation for TEP because it yields a linear program
amenable to efficient solution~\cite{frank2016or_opf}.  Molzahn and
Hiskens~\cite{molzahn2019opf_survey} provided a comprehensive survey of
relaxations and approximations, confirming that DC OPF remains the
workhorse for expansion planning despite advances in convex relaxations
of AC OPF.

\gtopt{} adopts the $B$-$\theta$ formulation (susceptance-angle) for
DC OPF, as recommended by Stott et al.~\cite{stott2009dc_opf}, with
extensions for resistive losses via piecewise-linear approximations and
phase-shifting transformers.

\subsection{Stochastic Dual Dynamic Programming}

The Stochastic Dual Dynamic Programming (SDDP) algorithm, introduced by
Pereira and Pinto~\cite{pereira1991sddp}, is the standard method for
solving large-scale multi-stage stochastic programs in hydrothermal
scheduling.  The algorithm constructs piecewise-linear approximations of
the expected cost-to-go function through iterative forward and backward
passes.

Shapiro~\cite{shapiro2011sddp_analysis} provided theoretical convergence
analysis of SDDP, and Philpott and de~Matos~\cite{philpott2013sddp_convergence}
extended the algorithm with risk-averse sampling.  Recent advances include
batch learning SDDP~\cite{leclercq2023batch_sddp}, which frames the
algorithm as Q-learning for improved convergence, and extensions to
stochastic OPF under uncertainty~\cite{fletcher2024sddp_opf}.

In the Julia ecosystem, Dowson and Kapelevich~\cite{dowson2020sddp_jl}
developed SDDP.jl, a modular SDDP implementation, together with a set
of reference problems used here for validation
(Section~\ref{sec:validation}).  \gtopt{} provides a C++ SDDP
implementation with backward-pass apertures, multi-cut and Markov-chain
variants, elastic-filter feasibility cuts, and boundary cuts that couple
studies of different horizons (Section~\ref{sec:multiscale}).

\subsection{Energy Storage in Expansion Planning}

The integration of battery energy storage systems (BESS) into expansion
planning has received significant attention.  Go et al.~\cite{Go2016ess_expansion}
assessed the economic value of co-optimized grid-scale storage investments.
Frazier et al.~\cite{frazier2023bess_milp} presented an integrated MILP
for generation, transmission, and storage expansion with Benders
decomposition.

Sandia National Laboratories released QuESt Planning~\cite{sandia2024quest},
an open-source capacity expansion tool focused on storage, while
IDAES-GTEP~\cite{idaes_gtep2024} provides a modular Pyomo-based framework
for GTEP.  \gtopt{} supports BESS with detailed charge/discharge efficiency
modeling, state-of-charge tracking across stages, and co-located
generation-storage coupling via its converter mechanism.

\subsection{Integrated Gas and Electricity Networks}
\label{sec:lit_gas}

Steady-state gas transmission is governed by the Weymouth equation,
which is linear in the squared pressure and quadratic in the
flow~\cite{menon2005gas}.  De~Wolf and Smeers~\cite{dewolf2000gas}
solved the resulting gas transmission problem with an extension of the
simplex algorithm, and Tomasgard et al.~\cite{tomasgard2007gas}
approximated the Weymouth cone by tangent planes inside an LP of the
gas value chain.  R\'ios-Mercado and
Borraz-S\'anchez~\cite{riosmercado2015review} reviewed the optimization
problems of gas transportation, and Borraz-S\'anchez et
al.~\cite{borraz2016relaxations} proposed convex relaxations for
expansion planning.  On the power side, Liu et al.~\cite{liu2009scuc_gas}
added gas transmission constraints to security-constrained unit
commitment, and Correa-Posada and
S\'anchez-Mart\'in~\cite{correaposada2015gas} integrated both networks
in short-term operation.  Ordoudis et al.~\cite{ordoudis2019gas} and
Ratha et al.~\cite{ratha2020affine} cleared coupled electricity and gas
markets with tangent-plane Weymouth constraints and line pack, Hu and
Ryan~\cite{hu2019stochastic} compared stochastic and deterministic
scheduling of a coupled system, and Bent et al.~\cite{bent2018joint}
planned joint transmission expansion with market feedbacks.

Saldarriaga~\cite{saldarriaga2013msc, saldarriaga2013holistic,
saldarriaga2018phd, saldarriaga2019stochastic} developed integrated
planning models for gas and electricity, from distribution networks to
multi-stage stochastic transmission planning; his doctoral thesis
derives the compressor-fuel model adopted in Section~\ref{sec:gas}.
Later work of the same group addressed seasonal storage in integrated
systems~\cite{martinez2022storage} and the LNG supply
chain~\cite{ibarra2019gnl, restrepo2021gnl}.  Open tools and data
support this line of research: GasModels.jl~\cite{gasmodels} and
GasPowerModels.jl~\cite{gaspowermodels} implement exact and relaxed gas
formulations in Julia, and GasLib~\cite{schmidt2017gaslib,
pfetsch2015nominations} provides gas network instances.

Most of these models are nonlinear or mixed-integer, or they treat the
gas system as a separate market.  \gtopt{} instead keeps the gas
network as a pure LP inside the same multi-stage hydrothermal model:
exact tangent envelopes on radial networks, a sequential LP for meshed
ones, and compressor fuel as Taylor planes in the squared pressure.
The gas and power duals are therefore consistent, and the network can
be solved by SDDP together with reservoirs and storage.

\subsection{Open-Source Power System Tools}

Table~\ref{tab:tool_comparison} compares \gtopt{} with the leading
open-source GTEP tools.  PyPSA~\cite{brown2018pypsa} is the most widely
adopted, offering multi-sector coupling and a large community, and
targets linear expansion and dispatch rather than stochastic
hydrothermal scheduling.
GenX~\cite{jenkins2017genx} provides a flexible Julia/JuMP framework but
lacks native DC OPF and hydro cascade modeling.  pandapower~\cite{thurner2018pandapower}
excels at power flow analysis but does not support expansion planning.

\begin{table*}[t]
  \centering
  \caption{Comparison of Open-Source GTEP and Power System Tools}
  \label{tab:tool_comparison}
  \begin{tabular}{@{}lcccccc@{}}
    \toprule
    \textbf{Feature} & \textbf{GTOpt} & \textbf{PyPSA} & \textbf{GenX}
      & \textbf{IDAES-GTEP} & \textbf{QuESt} & \textbf{PLP} \\
    \midrule
    Language          & C++26  & Python  & Julia   & Python  & Python  & Fortran \\
    License           & AGPL-3\textsuperscript{a} & MIT & GPL-2 & BSD-3 & BSD-3 & --\textsuperscript{b} \\
    DC OPF (Kirchhoff)& \checkmark & \checkmark & Transport & \checkmark & -- & \checkmark \\
    Hydro cascades    & \checkmark & Partial & -- & -- & -- & \checkmark \\
    Battery storage   & \checkmark & \checkmark & \checkmark & Partial & \checkmark & Partial \\
    SDDP solver       & \checkmark & -- & -- & -- & -- & \checkmark \\
    Capacity expansion& \checkmark & \checkmark & \checkmark & \checkmark & \checkmark & -- \\
    Multi-scenario    & \checkmark & Partial & Partial & -- & -- & \checkmark \\
    Multi-sector      & \checkmark & \checkmark & Partial & -- & -- & -- \\
    Line losses       & \checkmark & \checkmark & \checkmark & -- & -- & \checkmark \\
    Water rights      & \checkmark & -- & -- & -- & -- & \checkmark \\
    Unit commitment   & \checkmark & \checkmark & \checkmark & \checkmark & -- & -- \\
    Emission caps     & \checkmark & \checkmark & \checkmark & -- & -- & -- \\
    Gas network (pressure) & \checkmark & Transport & -- & -- & -- & -- \\
    Cut-coupled scales& \checkmark & -- & -- & -- & -- & Partial \\
    Parquet I/O       & \checkmark & -- & -- & -- & -- & -- \\
    Web GUI           & \checkmark & -- & -- & -- & \checkmark & -- \\
    \bottomrule
  \end{tabular}
  \\[3pt]
  {\footnotesize \textsuperscript{a}AGPL-3.0-or-later with a solver-plugin
   exception, or commercial.
   \textsuperscript{b}Maintained by the Chilean system operator (CEN).}
\end{table*}

\gtopt{} distinguishes itself through:
\begin{enumerate}
  \item one formulation used from daily operation to multi-decade
    planning, with the time scales coupled by SDDP cuts;
  \item complete hydro cascade modeling with reservoirs, waterways,
    turbines, pumps and seepage;
  \item a generic water-rights layer for irrigation agreements with
    volume-dependent dynamic bounds;
  \item three solution methods, including SDDP and cascade
    decomposition;
  \item an auditable model: a published formulation, exportable LP
    files and structural fingerprints;
  \item native Apache Parquet I/O and converters from the commercial
    models used in the region.
\end{enumerate}
